% Part: incompleteness
% Chapter: representability-in-q
% Section: comp-representable

\documentclass[../../../include/open-logic-section]{subfiles}

\begin{document}

\olfileid[pt]{inc}{req}{crq}
\olsection{Funções computáveis são representáveis em $\Th{Q}$}

\begin{thm}
Toda função computável é representável em~$\Th{Q}$.
\end{thm}

\begin{proof}
Para fixar as ideias, e usando a Tese de Church--Turing, digamos que
uma função é computável sse é recursiva geral. As funções recursivas
gerais são aquelas que podem ser definidas a partir da função
zero~$\Zero$, da função sucessor~$\Succ$ e da função de
projeção~$\Proj{n}{i}$ usando composição, recursão primitiva e
minimização regular. Por \olref[pri]{lem:prim-rec}, toda função~$h$ que
pode ser definida a partir de $f$ e~$g$ também pode ser definida por
composição e minimização regular a partir de $f$, $g$, $\Zero$,
$\Succ$, $\Proj{n}{i}$, $\Add$, $\Mult$, $\Char{=}$. Por conseguinte,
uma função é recursiva geral sse pode ser definida a partir de $\Zero$,
$\Succ$, $\Proj{n}{i}$, $\Add$, $\Mult$, $\Char{=}$ usando composição
e minimização regular.

Além disso, mostramos que as funções básicas em questão são
representáveis em~$\Th{Q}$
(\cref{inc:req:bre:prop:rep-zero,inc:req:bre:prop:rep-succ,inc:req:bre:prop:rep-proj,inc:req:bre:prop:rep-id,inc:req:bre:prop:rep-add,inc:req:bre:prop:rep-mult}),
e que toda função definida a partir de funções representáveis por
composição ou minimização regular
(\olref[cmp]{prop:rep-composition},
\olref[min]{prop:rep-minimization}) também é representável. Logo, toda
função recursiva geral é representável em~$\Th{Q}$.
\end{proof}

\begin{explain}
Mostramos que o conjunto das funções computáveis pode ser
caracterizado como o conjunto das funções representáveis em $\Th{Q}$.
Na verdade, a prova é mais geral. Pela definição de representabilidade,
não é difícil ver que toda teoria que estenda $\Th{Q}$ (ou na qual se
possa interpretar $\Th{Q}$) pode representar as funções computáveis.
Reciprocamente, porém, em todo sistema de !!{derivation} no qual a
noção de !!{derivation} seja computável, toda função representável é
computável. Assim, por exemplo, o conjunto das funções computáveis pode
ser caracterizado como o conjunto das funções representáveis na
aritmética de Peano ou até na teoria dos conjuntos de Zermelo--Fraenkel.
Como observou G\"odel, isso é um tanto surpreendente. Veremos que, no
que diz respeito à demonstrabilidade, as questões são muito sensíveis
à teoria considerada; aproximadamente, quanto mais fortes os axiomas,
mais se pode provar. Mas, em uma ampla classe de teorias axiomáticas,
as funções representáveis são exatamente as funções computáveis;
teorias mais fortes não representam mais funções, desde que sejam
axiomatizáveis.
\end{explain}

\end{document}
